\documentclass[border=10pt,crop=true]{standalone}
\usepackage{circuitikz}
% Shared style for 电子学一 Chapter 2 op-amp circuit figures (CircuiTikZ).
\usetikzlibrary{calc}

\ctikzset{
  bipoles/length=0.95cm,
  bipoles/thickness=1,
  resistors/scale=1,
  capacitors/scale=1,
  label distance=2.5mm,
}

\tikzset{
  opfig/.style={font=\small},
  opfigThin/.style={font=\scriptsize},
}

% Geometry (cm) — keep identical across all op-amp schematics
\def\OpInLen{1.55}    % label → input terminal
\def\OpInGap{0.08}    % terminal → wire to op amp +
\def\OpAmpWire{1.00}  % terminal side → op amp + anchor
\def\OpOutLen{1.05}   % op amp out → output terminal
\def\OpJuncL{0.50}    % v_- → feedback junction (left)
\def\OpFbDown{0.55}   % v_- straight down before R_1
\def\OpRvert{1.20}    % R_1 vertical span to ground
\def\OpInSep{0.55}    % dual-input spacing (v_1 / v_2)
\def\OpFbStub{0.40}   % follower: stub right of output before dropping
\def\OpFbClear{0.22}  % follower: below op-amp south for horizontal feedback


% 绝对坐标布局（cm，原点在运放几何中心附近）
\def\OpExOAx{1.05}          % 运放锚点 x（使 VM 落在 Xfb、O 落在 Xo）
\def\OpExXo{-0.65}          % 结点 O（须在 Xfb 右侧，避开 R_2 竖通道）
\def\OpExXfb{-1.50}         % 反馈节点
\def\OpExYbus{-2.50}        % R_4 水平母线（全图最低）
\def\OpExXoutR{2.00}       % 反馈上折返竖线 x
\def\OpExRhor{1.45}        % O/VM → R → 端子
\def\OpExRtwoDrop{1.55}     % O ↓ R_2 ↓ 地（止于 Ybus 之上）

\begin{document}
\begin{circuitikz}[american, scale=1.0, transform shape]
  \coordinate (OAC) at (\OpExOAx, 0);
  \node[op amp, noinv input up](OA) at (OAC) {};

  % 结点 O、反馈节点：与 +/- 引脚 y 对齐，x 取规定值
  \path let \p1=(OA.+), \p2=(OA.-) in
    coordinate (O) at (\OpExXo, \y1)
    coordinate (VM) at (\OpExXfb, \y2);

  % ===== 同相端（+）：+ ← O ← R_1 ← v_1；O ↓ R_2 ↓ 地 =====
  \draw (OA.+) -- (O);
  \node[circ, fill=black, inner sep=0.9pt] at (O) {};
  \node[opfigThin, above=0.55mm] at (O) {$O$};

  \coordinate (V1j) at ($(O)+(-\OpExRhor, 0)$);
  \draw (O) -- (V1j);
  \draw (V1j) to[R, l=$R_1$, l^=] ++(-0.82, 0)
    to[short, o-] ++(-\OpInGap, 0) node[left, opfig]{$v_1$};

  \draw (O) to[R, l_=$R_2$] ++(0, -\OpExRtwoDrop)
    node[ground, scale=0.85]{};

  % ===== 反相端（-）：- ← VM ← R_3 ← v_2（与上路平行）=====
  \draw (OA.-) -- (VM);
  \node[circ, fill=black, inner sep=0.9pt] at (VM) {};

  \coordinate (V2j) at ($(VM)+(-\OpExRhor, 0)$);
  \draw (VM) -- (V2j);
  \draw (V2j) to[R, l=$R_3$, l^=] ++(-0.82, 0)
    to[short, o-] ++(-\OpInGap, 0) node[left, opfig]{$v_2$};

  % ===== R_4：VM ↓(Xfb) → Ybus 右行 → XoutR ↑ → Output =====
  \coordinate (FBL) at (\OpExXfb, \OpExYbus);
  \coordinate (FBR) at (\OpExXoutR, \OpExYbus);
  \coordinate (OutRise) at (\OpExXoutR, 0);
  \draw (VM) -- (FBL);
  \draw (FBL) to[R, l_=$R_4$] (FBR);
  \draw (FBR) -- (OutRise |- OA.out) -- (OA.out);

  % ===== v_o =====
  \draw (OA.out) to[short, -o] ++(\OpOutLen, 0)
    node[right, opfig]{$v_o$};
\end{circuitikz}
\end{document}
